\documentclass{article}
\usepackage{times}
\usepackage{a4wide}
\usepackage{latexsym}
\usepackage{mathrsfs}
\usepackage{amsmath}
\usepackage{amssymb}
\usepackage{amsthm}
\usepackage{ifthen}
\usepackage{stmaryrd}
\usepackage{algorithm}
\usepackage{ctex}
\usepackage{algorithmic}
\usepackage{algpseudocode}
\usepackage{graphics}
\usepackage{epsfig}
\usepackage{setspace}
\setstretch{1}
%\usepackage{ramacros}

\addtolength{\textwidth}{20mm}
\addtolength{\oddsidemargin}{-10mm}
\addtolength{\textheight}{18mm}
\addtolength{\topmargin}{-9mm}


\newcounter{exercise}
\setcounter{exercise}{0}
\newcommand{\exercise}{
        \addtocounter{exercise}{1}
        \vspace{0.2in}
        \noindent
        {\bf \theexercise .}
        }


\newcommand{\solution}{
    \vspace{0.2in}
    \noindent
    {\bf Solution:}
    }


\newcommand{\REMARK}[1]{}


\newcommand{\NEWPART}{\vspace{.1in}
              \noindent}

\newcommand{\<}{
    \langle}

\renewcommand{\>}{
    \rangle}

\newcommand{\ceil}[1]{\left\lceil #1 \right\rceil}

\pagestyle{plain} \pagenumbering{arabic}

\title{{\bf Assignment 1} \\ {\large Deadline: Mar. 22}}

\author{125033910296 李博洋}
\date{}

\begin{document}
\maketitle

%{\large \noindent
%\begin{tabular}{lcl}
%{Jan 17 2007}.
%\end{tabular}
%}

{\large


\begin{exercise}
Prove the correctness of Dijkstra's shortest path algorithm.
\end{exercise}

\solution{
通过归纳法+反证法进行证明\\
\textbf{初始化：}\\
$dist[s] = 0$, $dist[s]$ 是到s的最短路。\\
对于每次引入的点 $u_i$，只要在该步已有 $dist[u_i]$ 为到 $u_i$ 的最短路时，仍能推出更新后的 $dist[u_{i+1}]$ 也是到 $u_{i+1}$ 的最短路，即可证明 Dijkstra 算法的正确性。\\
\textbf{反证：}\\
当$dist[u_i]$为最短路时，存在一条路$P$，使$L[P] < dist[u_{i+1}]$。
记算法迭代的过程中，$dist[u_i]$为最短路时，已加入最短路的点的集合为$S_i$。
由于$L[P] < dist[u_{i+1}]$，则$P$中至少有一个点不在$S_i$中，记第一个不在$S_i$中的点为$x$。
$P$上有边$(y, x), y \in S_i, x \notin S_i$，则$L[P_{s \to x}] \leq L[P] < dist[u_{i+1}]$。\\
而又因为$y \in S_i$，则$dist[y]$为最短路，有
$dist[x] \leq dist[y] + w(y, x) \leq L[P_{s \to y}] + w(y, x) = L[P_{s \to x}]$。\\
而又因在 Dijkstra 的每次迭代中，算法总是选取当前 $dist$ 最小的未确定点作为 $u_{i+1}$，因此有$dist[u_{i+1}] \leq dist[x]$。\\
由上可得$dist[u_{i+1}] \leq dist[x] \leq L[P_{s \to x}] < dist[u_{i+1}]$矛盾。\\
因此，Dijkstra 算法的正确性得证。
}


\begin{exercise}
In the 2SAT problem, you are given a set of clauses, where each clause is the disjunction of two literals. 
You are looking for a way to assign a value true or false to each of the variables so that all clauses are satisfied, 
that is, there is at least one true literal in each clause. Find a way of solving 2SAT efficiently.
[for detailed explanation, referred to 3.28 of DPV07]
\end{exercise}

\solution{
对于任何clause，可以如下表转换为两个推导式\\
\begin{table}[h]
    \centering
    \begin{tabular}{|c|c|c|}
        \hline
        Clause & 推导式  \\
        \hline
        $a \lor b$ & $a \to b$ 和$ \neg b \to \neg a$ \\
        \hline
        $a \lor \neg b$ & $a \to \neg b $和$ b \to a$\\
        \hline
        $\neg a \lor b$ & $a \to b$ 和$ \neg b \to \neg a$ \\
        \hline
        $\neg a \lor \neg b$ & $a \to \neg b$ 和 $b \to a$ \\
        \hline
    \end{tabular}
\end{table}\\
由此可以建图：对每个literal $x_i$，建两个点$x_i$和$\neg x_i$，对于每个clause，根据上表的推导式，建立对应的有向边。\\
由上述语义可知，同一SCC内的所有点，代表的条件要么都满足，要么都不满足。此时，若存在一组点$x_i$和$\neg x_i$在同一个SCC中，则说明该问题无解，反之，则说明有解。
因此可以通过kosaraju算法求出图中的所有的SCC，同时遍历所有点，检查同一SCC内的所有点是否存在$x_i$和$\neg x_i$，若存在，则说明该问题无解，反之，则说明有解。\\
当该问题有解时，可以用如下输出方案：\\
对每个SCC，选取其中任意一个点$x_j$，若该点所表示的 literal 在此之前未被赋值，则赋值为 true；否则沿用该 literal 已赋的值。然后从该点开始BFS该SCC中的所有点，
根据有向边的含义，将其所代表的literal赋值为true或false。
}


\begin{exercise}
Given $G$, calculate the minimum number of additional directed edges required to make the entire graph $G$ strongly connected.
\end{exercise}

\solution{
分步证明：\\
\textbf{1}: 首先使用kosaraju算法求出图$G$的强连通分量，并进行缩点，把所有SCC合并为一个点，得到一个DAG $D$。\\
\textbf{2}: 进行判断：\\
\textbf{2-1}: 如果DAG $D$中仅存在一个点，则说明此时图$G$强连通，不需要添加任何边。\\
\textbf{2-2}: 如果DAG $D$中存在多个点，由\textbf{引理：每一个DAG中至少有一个 source 和一个 sink} 可知，DAG $D$中至少存在一个source和一个sink，且每当增加一条从
sink到source的边，则可以使得DAG $D$ 减少一个source和一个sink。而当DAG $D$中仅存在source或sink时，增加从任何一个点到source的边或sink到任何一个点的边，则可
以使得DAG $D$ 减少一个source或一个sink。\\
同时，构造方案为：从每个sink点，向任何\textbf{从该source点不可达该sink点}且\textbf{未被其他sink点连边的点}的source点连边。
若不存在这样的点，则从该sink点向任意一个source点连边。若剩余多的source点，则从该source点向任意一个sink点连边。此时可保证任意两点可达。\\
因此，至少需要增加的边数为$\max(\mathrm{NUM}(\mathrm{source}), \mathrm{NUM}(\mathrm{sink}))$，其中$\mathrm{NUM}(\mathrm{source})$表示DAG $D$中source的数量，$\mathrm{NUM}(\mathrm{sink})$表示DAG $D$中sink的数量。
}

\begin{exercise}
A \emph{scorpion} is an undirected graph $G$ of the following form: there are three
special vertices, called the sting, the tail, and the body, of degree $1$, $2$,
and $n - 2$, respectively. The sting is connected only to the tail; the tail is
connected only to the sting and the body; and the body is connected to
all vertices except the sting. The other vertices of $G$ may be connected to
each other arbitrarily.

\centerline{\includegraphics{figures/scorpion}}


Give an algorithm that makes only $O(n)$ probes of the adjacency matrix of
$G$ and determines whether $G$ is a scorpion.
\end{exercise}

\solution{
算法的核心是利用蝎子图的\textbf{度数特征}快速定位关键顶点（body/tail/sting），避免遍历整个邻接矩阵（复杂度 $O(n^2)$），仅通过 $O(n)$ 次邻接矩阵探针完成判定：
\begin{description}
\item[1]: 定位 body 候选顶点
蝎子图的核心特征是body的度数为n-2（仅不与sting相连），而sting/tail的度数$\leq$2、普通顶点的度数介于2和n-2之间。
算法从任意顶点（如顶点$v_0$）开始，通过以下逻辑快速锁定body:
\begin{itemize}
    \item 如果顶点$d(v_i) = n-2$，则确定$v_i$为body，且其唯一非邻点为sting候选。
    \item 如果顶点$d(v_i) \leq 2$，说明$v_i$是sting/tail，跳转到其唯一/任一邻点（必连body）继续检查。
    \item 若 $d(v_i) > 2$ 且 $d(v_i) < n-2$，说明$v_i$是普通顶点，直接跳转到$v_i$的任意邻点重复此步骤（普通顶点必与body相连，因此最多1次跳转即可找到body）
\end{itemize}
\item[2]: 验证sting和tail的合法性
假设找到候选body为 $b$，执行以下验证（均为$O(n)$次探针）：
\begin{itemize}
    \item 找到唯一不与$b$ 相连的顶点 $s$（仅可能是 sting）
    \item 遍历$s$的邻接行（$O(n)$ 次探针），统计度数：若度数$\neq$1，直接判定非蝎子图；否则s的唯一邻点即为tail候选$t$；
    \item 遍历$t$的邻接行（$O(n)$ 次探针），统计度数并记录邻点：若度数$\neq$2，或邻点不是$s$和$b$，判定非蝎子图
\end{itemize}
\item
若以上验证均通过，则说明$G$是蝎子图，反之，则说明$G$不是蝎子图。
\end{description}

}


\begin{exercise}
You are the Chief Engineer of a galactic trade network. The network consists of $N$ star systems and $M$ one-way jump gates. Each jump gate from system $u$ to $v$ has a toll cost $w$.

Within this galaxy, some systems are part of ``Trade Federations''. A Trade Federation is a maximal group of systems where every system can reach every other system using existing jump gates.

\begin{description}
\item[Internal Travel]: Traveling between two systems within the same federation is subsidized—the toll cost is effectively
\item[External Travel]: Traveling between two different federations requires paying the toll of the jump gate used.
\end{description}

The Galactic Council wants to build new bidirectional bridges between different Trade Federations to ensure that every system in the galaxy can reach every other system.

However, they want to spend the minimum amount possible on the tolls of the jump gates that will ''connect'' these federations.

Find the minimum total toll cost of a set of existing jump gates that, if converted into bidirectional bridges, would make the entire galaxy strongly connected.
\end{exercise}

\solution{
\begin{description}
    \item[步骤一]: 首先使用kosaraju算法求出图$G$的强连通分量，并进行缩点，把所有SCC合并为一个点，得到一个DAG $D$。
    \item[步骤二]: 然后使用kruskal算法求出DAG $D$的最小生成树，并输出最小生成树的边权和。
\end{description}
因此，最小总通行费用为DAG $D$的最小生成树的边权和
}


\begin{exercise}
 \textbf{Minimum Bottleneck Spanning Tree}: Given a connected graph $ G $ with positive edge costs, find a
 spanning tree that {minimizes the most expensive edge}.
\end{exercise}

\solution{
\textbf{首先，证明最小生成树一定是最小瓶颈生成树：}\\
\textbf{反证：}\\
设最小生成树中的最大边权为$w$, 如果最小生成树不是瓶颈生成树的话，则瓶颈生成树的所有边权都小于$w$，只需删去原最小生成树中的最长边，
用瓶颈生成树中的一条边来连接删去边后形成的两棵树，得到的新生成树一定比原最小生成树的权值和还要小，这样就产生了矛盾．因此，最小生成树一定是最小瓶颈生成树。\\
\textbf{然后，通过Kruskal算法求解最小生成树：}
Kruskal算法是一种贪心算法，每次选择当前权值最小的边，同时检查每次加入的边是否会导致环的产生，若加入后形成环，则不加入该边。反之则加入。
由于最小生成树一定是最小瓶颈生成树，因此Kruskal算法求解的结果一定是最小瓶颈生成树。\\
\begin{algorithm}[H]
    \caption{Kruskal's Algorithm for Minimum Bottleneck Spanning Tree (MBST)}
    \label{alg:kruskal}
    \begin{algorithmic}[1] 
        \Require 
            $G = (V, E)$: 带权连通无向图，$V$ 为顶点集，$E$ 为边集；
            $w(e)$: 边 $e$ 的权值（正整数）
        \Ensure 
            $T$: 最小瓶颈生成树（MBST）；
            $b$: MBST 的瓶颈值（树中最大边的权值）
        
        \State \textbf{初始化} $T \leftarrow \emptyset$, $b \leftarrow 0$  \Comment{初始化生成树和瓶颈值}
        \State 将边集 $E$ 按权值 $w(e)$ \textbf{升序}排序，记为 $E_{\text{sorted}}$  \Comment{预处理步骤}
        \State 初始化并查集，每个顶点自成一个集合  \Comment{用于检测环}
        
        \For{每一条边 $e = (u, v) \in E_{\text{sorted}}$}
            \If{$u$ 和 $v$ 属于不同的集合}  \Comment{加入边不会形成环}
                \State 将边 $e$ 加入 $T$  \Comment{扩展生成树}
                \State $b \leftarrow w(e)$  \Comment{更新当前最大边权（瓶颈值）}
                \State 合并 $u$ 和 $v$ 所在的集合  \Comment{维护连通性}
                \If{$|T| = |V| - 1$}  \Comment{生成树已完成（n-1条边）}
                    \State \textbf{break}  \Comment{提前终止循环，优化效率}
                \EndIf
            \EndIf
        \EndFor
        
        \State \Return $T, b$
    \end{algorithmic}
\end{algorithm}

\begin{algorithm}[H]
    \caption{Union-Find Operations}
    \label{alg:uf}
    \begin{algorithmic}[1]
        \Function{Find}{$x$}  \Comment{查找顶点x的根节点（路径压缩）}
            \If{$parent[x] \neq x$}
                \State $parent[x] \leftarrow \text{Find}(parent[x])$  \Comment{路径压缩，降低后续查找成本}
            \EndIf
            \State \Return $parent[x]$
        \EndFunction
        
        \Function{Union}{$x, y$}  \Comment{合并x和y所在的集合（按秩合并）}
            \State $x_{\text{root}} \leftarrow \text{Find}(x)$
            \State $y_{\text{root}} \leftarrow \text{Find}(y)$
            \If{$x_{\text{root}} = y_{\text{root}}$}
                \State \Return $\text{false}$  \Comment{已在同一集合，合并失败}
            \EndIf
            \If{$rank[x_{\text{root}}] < rank[y_{\text{root}}]$}
                \State $parent[x_{\text{root}}] \leftarrow y_{\text{root}}$
            \Else
                \State $parent[y_{\text{root}}] \leftarrow x_{\text{root}}$
                \If{$rank[x_{\text{root}}] = rank[y_{\text{root}}]$}
                    \State $rank[x_{\text{root}}] \leftarrow rank[x_{\text{root}}] + 1$
                \EndIf
            \EndIf
            \State \Return $\text{true}$  \Comment{合并成功}
        \EndFunction
    \end{algorithmic}
\end{algorithm}

}

\end{document}
